\begin{abstract}
    随着全球能源结构转型和``双碳''目标的推进，分布式储能系统在电力系统中的应用日益广泛。本项目开发计划书详细阐述了``智能能源管理与数据科学平台''的开发规划，该平台旨在通过集成机器学习预测与数学优化求解技术，实现储能系统的智能化、精细化管理。

    本文档遵循IEEE 1058软件项目管理计划标准，从项目概述、组织管理、技术实施、质量保障四个维度进行了系统性规划。项目采用``预测+优化''的两阶段技术路线，利用机器学习算法实现负载预测，利用Gurobi混合整数规划求解器实现储能调度优化。平台基于Flask后端与Flutter Web前端构建，部署于Google Cloud Platform，实现云原生架构的弹性伸缩与高可用保障。

    通过本计划书的实施，预期可实现电费节省15\%--25\%的经济目标，同时为储能系统运营商提供一套可靠、透明、易用的智能决策支持工具。

    \thusetup{
        keywords = {储能优化, 负载预测, 混合整数规划, 机器学习, 云原生架构},
    }
\end{abstract}

\begin{abstract*}
    With the global energy transition and the advancement of ``carbon peaking and carbon neutrality'' goals, distributed energy storage systems are increasingly applied in power systems. This Project Development Plan elaborates on the development of the ``Smart Energy Management \& Data Science Platform'', which aims to achieve intelligent and refined management of energy storage systems through the integration of machine learning prediction and mathematical optimization techniques.

    This document follows the IEEE 1058 Software Project Management Plan standard, providing systematic planning from four dimensions: project overview, organization management, technical implementation, and quality assurance. The platform adopts a ``Predict + Optimize'' two-stage approach, utilizing machine learning algorithms for load forecasting and the Gurobi mixed-integer programming solver for optimal energy storage scheduling. The platform is built on a Flask backend and Flutter Web frontend, deployed on Google Cloud Platform, achieving elastic scaling and high availability through cloud-native architecture.

    Through the implementation of this plan, the expected economic goal of 15\%--25\% electricity cost savings can be achieved, while providing energy storage operators with a reliable, transparent, and user-friendly intelligent decision support tool.

    \thusetup{
        keywords* = {Energy Storage Optimization, Load Forecasting, Mixed-Integer Programming, Machine Learning, Cloud-Native Architecture},
    }
\end{abstract*}